\section*{第 \ref{ch:b2:wilson-cycle} 章——古老地块与威尔逊旋回}

\begin{solution}{exo:b2:wilson-cycle:1}
(a) $600 - 400 = \qty{200}{Ma}$，自边缘裂开至最后一次碰撞。
(b) $480 - 400 = \qty{80}{Ma}$。(c) 加里东造山作用是一系列分布在
约 2 亿年间的阶段，各阶段在不同地点开始的时间并不相同：这个词指明的是一个洋及其\omterm{def:b2:alpine-orogen:orogen}{造山带}，
而非某一时刻。年龄由地层年表给出，
或以百万年计。
\end{solution}

\begin{solution}{exo:b2:wilson-cycle:2}
(c) 东非裂谷，裂开；(a) 红海，年轻的洋；(d) 大西洋，
一个具有被动大陆边缘的宽阔洋盆；(e) 南美洲太平洋边缘，俯冲与洋盆闭合；
(b) 喜马拉雅，碰撞；
(f) 布列塔尼，一个已被夷平的\omterm{def:b2:alpine-orogen:orogen}{造山带}，此后裂谷可能再度开始。
\end{solution}

\begin{solution}{exo:b2:wilson-cycle:3}
(a) 在早志留世地层沉积之后，即 \qty{443.1}{Ma} 之后，
且在晚泥盆世地层沉积之前，即 \qty{358.9}{Ma} 之前，也就是泥盆纪末。
(b) 少于 $443.1 - 358.9 = \qty{84.2}{Ma}$；实际上更少，因为下部地层比志留纪底界更年轻，
上部地层比其顶界更古老。(c) 大洋中灰色浊积岩的沉积；
加里东碰撞中的褶皱与劈理；隆升并侵蚀，形成起伏不平的陆地表面；
红色角砾岩与砾岩的沉积；
瓦里斯坎造山作用中的缓倾斜；今日切割悬崖的侵蚀。
\end{solution}

\begin{solution}{exo:b2:wilson-cycle:4}
(a) 两者皆是：一个陆壳地块，穿越艾珀特斯洋时与大陆被洋盆隔开，
拼贴上来之后，便是一个以断层为界、历史与其邻区不同的地块。
(b) 两者皆非：它是一个大陆，而非附加于另一大陆之上的地块。(c) \omterm{def:b2:wilson-cycle:terrane}{地体}，而非\omterm{def:b2:wilson-cycle:terrane}{微大陆}：
其地壳属于火山弧，而非大陆的碎片。(d) 两者皆非：它是\omterm{def:b2:wilson-cycle:terrane}{地体}之间的边界，
即闭合洋盆的遗迹。
\end{solution}

\begin{solution}{exo:b2:wilson-cycle:5}
$3000 - 2600 = 400$，$2600 - 1500 = 1100$，$1500 - 875 = 625$，$875 - 600 = 275$，$600 -
262.5 = 337.5$；平均值 $2737.5/5 = \qty{547.5}{Ma}$。(b) $1100/275 = 4.0$。(c) $(1500 -
262.5)/3 = \qty{412.5}{Ma}$。剔除可疑数据点后，平均值变化了四分之一，
而间隔变化达四倍：这只是粗略的复现，而非周期。
\end{solution}

\begin{solution}{exo:b2:wilson-cycle:6}
(a) $3000/(70 - 25) = \qty{66.7}{Ma}$。(b) $3000/70 = \qty{42.9}{Ma}$。(c) 洋脊一旦俯冲，
剩余的洋底便属于远端板块，该板块携带着远端大陆；
闭合速率即为该板块相对于海沟的速度。
\end{solution}

\begin{solution}{exo:b2:wilson-cycle:7}
(a) $700/8 = \qty{87.5}{km}$；$700/30 = \qty{23.3}{km}$。(b) 瓦里斯坎带：
按每千米 \qty{30}{\celsius} 的梯度，中地壳会熔融；按每千米 \qty{8}{\celsius} 的梯度，
熔点则位于即便增厚地壳底部之下。(c) 低梯度描述的是在俯冲带中
被迅速带入深部、来不及升温的岩石。停留在增厚地壳
深部的岩石会在数百万年间向正常地温线升温，并
随后在较浅深度达到熔融。
\end{solution}

\begin{solution}{exo:b2:wilson-cycle:8}
(a) $\tau\ln(3/0.3) = 26.3\times2.303 = \qty{60.6}{Ma}$。(b) $290 - 60.6 = \qty{229}{Ma}$，位于
晚三叠世。(c) $\tau = (290 - 251.9)/\ln10 = \qty{16.5}{Ma}$。(d) 增厚地壳的伸展垮塌，
使其减薄，并直接降低地表高度；
此外，侵蚀速率更高、气候更湿润、岩石更软或起始海拔较低，也会导致这一结果。
\end{solution}

\begin{solution}{exo:b2:wilson-cycle:9}
(a) 缝合带，即艾珀特斯洋的遗迹，位于大西洋东侧的不列颠与爱尔兰陆地上；
来自艾珀特斯洋对岸的劳伦大陆地壳，位于欧洲一侧的罗克尔岩附近。
新裂谷切割于老缝合带以西：它是在缝合带附近张开的，而非沿缝合带张开。
(b) 缝合带以北是劳伦大陆边缘的基底与盖层，以南则是阿瓦隆尼亚的基底与盖层，
两者历史不同，且在早古生代的动物群也不同。(c) 裂谷出现于 8000 万年之后，
超过了冷却带可能变得坚固所需的 3000 万至 5000 万年；然而中大西洋确实沿阿巴拉契亚山脉发育。
\omterm{def:b2:alpine-orogen:orogen}{造山带}的热量会逐渐消退，但其断层、组构与成分则会保留，
并仍可引导裂谷的发育。
\end{solution}

\begin{solution}{exo:b2:wilson-cycle:10}
(a) $70 - 40/2 = \qty{50}{mm}$ 每年。(b) $\qty{1500}{km}/\qty{50}{km.Ma^{-1}} =
\qty{30}{Ma}$。(c) 洋盆以 $70 - 40 = \qty{30}{mm}$ 每年收窄：$3000 - 30\times30 =
\qty{2100}{km}$。(d) 年轻、炽热的岩石圈随之抵达海沟，
不再有古老、冷的板片在弧下脱水，弧\omterm{def:b2:magma-emplacement:magma}{岩浆}作用随之改变或停止。
\end{solution}

\begin{solution}{exo:b2:wilson-cycle:11}
(a) $h = \Delta T(\rho_m - \rho_c)/\rho_m = 25\times600/3300 = \qty{4.5}{km}$。(b) 额外的 \qty{25}{km} 全部，
相当于山脉高度的 5.5 倍。(c) $2700\times9.8\times25\,000 =
\qty{6.6e8}{Pa}$，约 \qty{0.66}{GPa}。(d) 在约 \qty{0.7}{GPa} 的压力以及增厚地壳根部的温度（数百度）下，
形成角闪岩相或麻粒岩相，
即古老地块中的片麻岩与\omterm{def:b2:crustal-melting:migmatite}{混合岩}。
\end{solution}

\begin{solution}{exo:b2:wilson-cycle:12}
(a) $S = 30\times10 = \qty{300}{km}$；$W_0 = 300 + 300 = \qty{600}{km}$；$T = 35\times600/300 =
\qty{70}{km}$。(b) $h = 35\times600/3300 = \qty{6.36}{km}$；$r = 4.5h = \qty{28.6}{km}$。
(c) $\Delta F = 2700\times9.8\times6364\times(35\,000 + 3182 + 14\,318) =
\qty{8.8e12}{N.m^{-1}}$，为西藏的 1.4 倍。(d) 这样的\omterm{def:b2:alpine-orogen:orogen}{造山带}比已在扩展的西藏更重：
发生\omterm{def:b2:himalaya-andes:collapse}{重力垮塌}，表现为高区伸展以及边缘逆冲。
\end{solution}

\begin{solution}{pb:b2:wilson-cycle:1}
\textbf{1.} $720/(540 - 500) = \qty{18}{km}$ 每百万年，即 \qty{18}{mm} 每年。
\textbf{2.} \qty{36}{mm} 每年：慢速扩张洋脊（20 至 \qty{55}{mm} 每年）。
\textbf{3.} $36\times(580 - 540) = \qty{1440}{km}$。
\textbf{4.} $36\times(580 - 500) = \qty{2880}{km}$。
\textbf{5.} 自 \qty{580}{Ma} 时宽度为零起，以恒定速率对称扩张，
且在 \qty{500}{Ma} 之前没有俯冲。
\textbf{6.} \qty{80}{Ma}，位于洋盆两侧，紧邻大陆边缘。它是
该洋盆中最冷、最厚、最致密的大洋岩石圈，最易下沉。
\textbf{7.} $64 - 36 = \qty{28}{mm}$ 每年。
\textbf{8.} $2880/28 = \qty{102.9}{Ma}$，即 \qty{500}{Ma} 之后：位于 \qty{397.1}{Ma}。
\textbf{9.} 洋脊距海沟 \qty{1440}{km}，以 $64 - 18 =
\qty{46}{mm}$ 每年接近海沟：经过 $1440/46 = \qty{31.3}{Ma}$，位于 \qty{468.7}{Ma}。
\textbf{10.} $2880 - 28\times31.3 = \qty{2004}{km}$。
\textbf{11.} $2004/40 = \qty{50.1}{Ma}$ 之后：位于 $468.7 - 50.1 = \qty{418.6}{Ma}$。
\textbf{12.} 位于 425 至 \qty{400}{Ma} 的斯堪的纳维亚碰撞期间。
\textbf{13.} $40\times8 = \qty{320}{km}$。
\textbf{14.} $W_0 = 400 + 320 = \qty{720}{km}$；$T = 35\times720/400 = \qty{63.0}{km}$。
\textbf{15.} $h = 28\times600/3300 = \qty{5.09}{km}$；$r = 4.5\times5.09 = \qty{22.9}{km}$。
\textbf{16.} $\Delta F = 2700\times9.8\times5091\times(35\,000 + 2545 + 11\,455) =
\qty{6.60e12}{N.m^{-1}}$，约比西藏多 \qty{8}{\percent}。
\textbf{17.} 要停止隆升并向外扩展：高区内部出现伸展与正断层，
边缘出现逆冲，正如今日的西藏。
\textbf{18.} $f = 600/3300 = 0.182$；$k = 2.303/11 = 0.209$ 每百万年；$\tau = 1/(kf)
= \qty{26.3}{Ma}$。
\textbf{19.} $3.0\,e^{-30/26.3} = \qty{0.96}{km}$。
\textbf{20.} $26.3\ln(3.0/0.5) = \qty{47.1}{Ma}$，即 \qty{410}{Ma} 之后：位于 \qty{362.9}{Ma}。
\textbf{21.} 多余的地壳为 $h/f$：$3.0/0.182 = \qty{16.5}{km}$ 降至 $0.5/0.182 =
\qty{2.75}{km}$，故移除了 \qty{13.75}{km} 的岩石。
\textbf{22.} $2700\times9.8\times13\,750 = \qty{3.6e8}{Pa}$，约为 \qty{0.36}{GPa}。
\textbf{23.} 是的：到晚泥盆世，模型\omterm{def:b2:alpine-orogen:orogen}{造山带}的平均海拔已降至约 \qty{1}{km}，
且仍在下降，其谷与脊之间仍保留起伏；红色砾石可在
\qty{358.9}{Ma} 之前很久便在其低洼处堆积，如锡卡尔角，那里的红层覆盖在一个
具有起伏的表面之上。
\textbf{24.} \textbf{\qty[text-series-to-math=true]{2880}{km}}，位于 \qty{500}{Ma}，
即俯冲开始之时。
\end{solution}
